%% Style by Bram Gotink, as a clone of the LyX style orginally used, but in clean code.
\documentclass[dutch]{article}
\usepackage[T1]{fontenc}
\usepackage[utf8]{inputenc}
% XeLaTeX/Tectonic negeren inputenc: zet deze tekens om, anders verdwijnen ze
% of worden ze fout gedrukt (§ werd ğ). Onder pdflatex verandert er niets.
\usepackage{iftex}
\ifTUTeX
  \usepackage{newunicodechar}
  \newunicodechar{§}{\S}
  \newunicodechar{‘}{\textquoteleft}
  \newunicodechar{’}{\textquoteright}
  \newunicodechar{“}{\textquotedblleft}
  \newunicodechar{”}{\textquotedblright}
\fi
\usepackage{geometry}
\geometry{verbose,letterpaper,tmargin=2.5cm,bmargin=2.5cm,lmargin=2.5cm,rmargin=2.5cm}
\setlength{\parskip}{\medskipamount}
\setlength{\parindent}{0pt}
\usepackage{textcomp}
\usepackage{titlesec}
\usepackage{enumitem}
\usepackage{color}
\usepackage{ulem}
\usepackage{hyperref}
%\hypersetup{colorlinks=false,urlcolor=cyan,pdfborder=0 0 0}

% VTK Styling

% subsections worden gebruikt als artikels in het HR en Verkiezingsreglement,
% Deze nummering moet doorlopen
\counterwithout{subsection}{section}

% parts zijn "Titel I - De vereniging"
\def\thepart{\Roman{part}}
\def\labelpart{Titel}
\titleformat{\part}
            [hang]
            {\normalfont\huge\bfseries}
            {\labelpart~\thepart~-}
            {10pt}
            {}
% sections zijn "Artikel 1. Naam" enzo
\def\thesection{\arabic{section}}
\def\labelsection{Artikel}
\titleformat{\section}
            [hang]
            {\normalfont\Large\bfseries}
            {\labelsection~\thesection.}
            {10pt}
            {}
% subsections zijn "Artikel 1. Naam", enkel gebruikt in HR
% (daar zijn sections "Afdeling 1 - Naam", zie de .tex van HR)
\def\thesubsection{\arabic{subsection}}
\def\labelsubsection{Artikel}
\titleformat{\subsection}
			[hang]
			{\normalfont\large\bfseries}
			{\labelsubsection~\thesubsection.}
			{10pt}
			{}
% hoogste niveau lijsten zijn §1. etc
\def\theenumi{\S\arabic{enumi}}
\def\labelenumi{\theenumi.}
% tweede niveau wordt eigenlijke eerste niveau
\def\theenumii{\arabic{enumii}}
\def\labelenumii{\theenumii.}
% eerste niveau itemize komt enkel voor in een lijst
\def\labelitemi{}
% layout de lijsten wat mooier
\setlist[enumerate,1]{align=left,labelindent=0pt,labelsep=3em,leftmargin=5em}
\setlist[enumerate,2]{align=left}
\setlist[itemize,1]{align=left,labelindent=0pt,labelsep=3em,leftmargin=5em,labelwidth=0pt}

% locate changes and color them red.
\newenvironment{aanpassing}{\color{red}}{}
\newcommand{\aangepast}[1]{\begin{aanpassing}#1\end{aanpassing}}
\newcommand{\verwijderd}[1]{\aangepast{\sout{#1}}}
% nieuwe (voorgestelde) tekst in het groen
\definecolor{vtkgroen}{rgb}{0,0.55,0}
\newcommand{\toegevoegd}[1]{{\color{vtkgroen}#1}}
\newcommand{\draft}{\textit{\textbf{draft}}}

\usepackage{babel}



\begin{document}


\title{Gecoördineerde statuten Vlaamse Technische Kring vzw}
\date{1 juli 2024}

\maketitle


\part{De vereniging}

    \section{Rechtsvorm}
    \begin{enumerate}
    \item De vereniging wordt opgericht als een entiteit met rechtspersoonlijkheid, meer in het bijzonder een vereniging zonder winstoogmerk (hierna "vwz"\ genaamd) op grond van de wet van 27 juni 1921 (B.S. 1 juli 1921) betreffende verenigingen zonder winstoogmerk, de internationale verenigingen zonder winstoogmerk en de stichting, zoals gewijzigd bij wet d.d. 2 mei 2002, 16 januari 2003, 22 december 2003, 9 juli 2004, 27 december 2004, 24 maart 2019, 29 april 2019, 12 mei 2019 en 28 april 2020. De statuten worden aangepast conform de bepaling van het Wetboek van Vennootschappen en Verenigingen (hierna: het "WVV"\ genoemd).
    \end{enumerate}
    
    \section{Naam}
    \begin{enumerate}
    \item De vereniging draagt de naam "Vlaamse Technische Kring", afgekort "VTK".
    \end{enumerate}
    
    \section{Zetel}
    \begin{enumerate}
    \item De zetel van de vereniging is gevestigd te Studentenwijk Arenberg 6/1, 3001 Heverlee, gelegen in het Vlaamse Gewest.
    
    \item Het Bestuursorgaan heeft de bevoegdheid om de zetel te verplaatsen naar iedere plaats binnen het Nederlandse taalgebied van België en de nodige openbaarmakingsvereisten te vervullen. De zetelwijziging moet worden bekrachtigd in de statuten door de Algemene Vergadering op de eerstvolgende vergadering.
    \end{enumerate}
    
    \section{Doel en voorwerp}
    \begin{enumerate}
    \item De vereniging heeft tot belangeloos doel, als faculteitskring van de Faculteit Ingenieurswetenschappen van de KU Leuven, een aanvulling te verschaffen op de universitaire vorming van de ingenieursstudent bij het opgelegde programma van de Universiteit, in het bijzonder op sociaal, cultureel en economisch gebied. De vzw keert op straffe van nietigheid, rechtstreeks noch onrechtstreeks enig vermogensvoordeel uit aan de oprichters, de leden, de bestuurders of enig ander persoon, behalve in dit laatste geval, voor het hierboven bepaald belangeloos doel.
    \item De vereniging heeft tot maatschappelijk voorwerp het bevorderen van de werkzaamheid die de vriendschapsbanden onder haar leden en met andere studenten of afgestudeerden van de Faculteit Ingenieurswetenschappen versterken. Verder dient zij (op onderwijskundig vlak) het behartigen van de belangen van de studenten aan de Faculteit Ingenieurswetenschappen te verzekeren. Ze poogt eveneens de studenten aan deze faculteit in contact te brengen met het bedrijfsleven en voor hen ontspannende activiteiten te voorzien.
    \item De vereniging mag alle rechtshandelingen stellen die nodig of nuttig zijn voor het verwezenlijken van het doel, en mag daartoe alle nodige roerende en onroerende goederen in huur of gebruik nemen, verwerven, bezitten en vervreemden. Zij mag ook haar medewerking verlenen aan alle verenigingen, kringen of instellingen die een gelijkaardig doel als het hare nastreven.
    \item Daarnaast kan de vereniging alle activiteiten ontplooien die rechtstreeks of onrechtstreeks bijdragen tot de verwezenlijking van de voormelde ideële niet-winstgevende doelstellingen, met inbegrip van bijkomstige commerciële en winstgevende activiteiten binnen de grenzen van wat wettelijk is toegelaten en waarvan de opbrengsten te allen tijde volledig zullen worden bestemd voor de verwezenlijking van de ideële niet-winstgevende doelstellingen. 
    \end{enumerate} 
    
    \section{Duur}
    \begin{enumerate}
    \item De vereniging is opgericht voor onbepaalde duur. Zij kan te allen tijde ontbonden worden door de Algemene Vergadering, conform artikel \ref{ArtikelOntbinding}.
    \end{enumerate}
    
    % \section{Vermeldingen}
    % \begin{enumerate}
    % \item Alle akten, facturen, aankondigingen, bekendmakingen, brieven, orders, websites en andere stukken, al dan niet in elektronische vorm, uitgaande van de vereniging moeten de volgende gegevens vermelden:
    % \begin{itemize}
    %     \item de naam van de vereniging;
    %     \item de rechtsvorm, voluit of afgekort;
    %     \item de nauwkeurige aanduiding van de zetel van de vereniging;
    %     \item het ondernemingsnummer;
    %     \item het woord "rechtspersonenregister"\ of de afkorting "RPR", gevolgd door de vermelding van de rechtbank van de zetel van de vereniging;
    %     \item het e-mailadres en de website van de vereniging;
    %     \item in voorkomend geval, het feit dat de vereniging in vereffening is.
    % \end{itemize}
    % \end{enumerate}


\part{Leden}

     \section{Werkende leden}\label{ArtikelWerkendeLeden}
    \begin{enumerate}
    \item Het aantal werkende leden is onbeperkt, maar moet ten minste twee bedragen. De werkende leden zijn titularis van alle rechten en plichten die in de WVV-wet aan "leden"\ worden toegekend.\label{ParMinAantalLeden}
    % \item Zowel natuurlijke personen als rechtspersonen kunnen werkend lid worden van de vereniging. Indien het een rechtspersoon betreft benoemt deze een natuurlijke persoon als vaste vertegenwoordiger die wordt belast met de uitvoering van het mandaat in naam en voor rekening van de rechtspersoon.
    \item Natuurlijke personen moeten, om lid te kunnen zijn van de vereniging, ingedeeld zijn bij VTK, in overeenstemming met het Participatiereglement van de Studentenraad KU Leuven en indien van toepassing het Participatiereglement van LOKO.
    \label{ParWerkendLidMoetStudentZijn}
    \item Bovendien moeten zij volgens diezelfde reglementen stemgerechtigd zijn. \label{ParWerkendLidMoetStemgerechtigdZijn}
    \item Natuurlijke personen die aan de voorwaarden zoals beschreven in artikel \ref{ArtikelWerkendeLeden} \ref{ParWerkendLidMoetStudentZijn} en \ref{ParWerkendLidMoetStemgerechtigdZijn} voldoen worden verder "studenten aan de Faculteit Ingenieurswetenschappen"\ genoemd.
    \item De Algemene Vergadering kan beslissen om natuurlijke personen die niet voldoen aan artikel \ref{ArtikelWerkendeLeden}\ref{ParWerkendLidMoetStudentZijn}, maar wel aan \ref{ArtikelWerkendeLeden} \ref{ParWerkendLidMoetStemgerechtigdZijn} toe te laten als werkend lid van de vereniging.
    \item De natuurlijke personen die in een bepaald werkjaar als praesidiumlid werden verkozen naar aanleiding van de praesidiumverkiezingen zoals geregeld in het Huishoudelijk Reglement, worden op de eerstvolgende 15 juli volgend op de dag van de verkiezingen van rechtswege werkend lid van de vereniging, na bekrachtiging van de geldigheid van hun verkiezing door het Bestuursorgaan. Hun lidmaatschap geldt voor één jaar en is onbeperkt hernieuwbaar.
    \item Onverminderd de bepaling van artikel \ref{ArtikelWerkendeLeden} \ref{ParMinAantalLeden}, is enkel de Algemene Vergadering bevoegd om personen als werkend lid tot de vereniging toe te laten. De Algemene Vergadering bepaalt soeverein aan wie ze het lidmaatschap aanbiedt. Het lidmaatschap eindigt de 14de juli, volgend op de datum dat het lidmaatschap aanvang nam.
    \item Zoals gespecifieerd in artikel \ref{ArtikelAudit} van deze statuten verkiest de Algemene Vergadering ten minste 3 en hoogstens 5 personen van het Auditcomité als werkend lid. In afwijking van artikel \ref{ArtikelWerkendeLeden} \ref{ParWerkendLidMoetStudentZijn} en \ref{ParWerkendLidMoetStemgerechtigdZijn} dienen deze werkende leden geen student van de Faculteit Ingenieurswetenschappen te zijn.
    \item Alle beslissingen betreffende toetreding, uittreding of uitsluiting van de werkende leden worden door het Bestuursorgaan in een register ingeschreven binnen de 8 dagen na kennisname van de inwerkingtreding van de beslissing.
    \item Het Bestuursorgaan kan aan de werkende leden specifieke volmachten verlenen, eigen aan hun taak.
    \item Van de werkende leden worden geen financiële bijdragen of stortingen gevraagd.
    \end{enumerate}
    
    \section{Toegetreden leden}
    \begin{enumerate}
    \item Naast de werkende leden kent de vereniging drie categorieën van toegetreden leden: ``gewone leden'', ``ereleden'' en ``steunende leden''.
    \item De toegetreden leden hebben enkel de rechten en plichten die in deze statuten staan beschreven.
    \item Elke student aan de Faculteit Ingenieurswetenschappen kan zich kandidaat stellen als gewoon lid. Gewone leden betalen een jaarlijkse bijdrage waarvan het bedrag voor alle gewone leden gelijk is en door het Bestuursorgaan wordt bepaald bij het begin van elk werkingsjaar en die maximum EUR 20,- zal bedragen. Deze bijdrage geldt voor één werkjaar.
    \item Iedereen die niet studeert aan een Belgische universiteit of hogeschool kan zich kandidaat stellen als erelid. Ereleden betalen een jaarlijkse bijdrage waarvan het bedrag vrij gekozen wordt, maar met een minimumbedrag dat voor alle ereleden gelijk is en door het Bestuursorgaan wordt bepaald bij het begin van elk werkjaar en dat maximum EUR 100,- zal bedragen. Deze bijdrage geldt voor één jaar. Het Bestuursorgaan heeft zeggenschap over het al dan niet toelaten van ereleden en kan verkiezen het lidmaatschap aan bepaalde personen gratis aan te bieden. Het Bestuursorgaan kan beslissen om een student van een Belgische universiteit of hogeschool bij uitzondering alsnog toe te laten als erelid.
    \item Elke student van een Belgische universiteit of hogeschool die de doelstelling van de vereniging onderschrijft kan zich kandidaat stellen als steunend lid. Steunende leden betalen een jaarlijkse bijdrage waarvan het bedrag voor alle steunende leden gelijk is en door het Bestuursorgaan wordt bepaald bij het begin van elk werkjaar en die maximum EUR 100,- zal bedragen. Deze bijdrage geldt voor één jaar. Het Bestuursorgaan kan personen weigeren om toe te treden als steunend lid.
    % To Do: hier fixen
    \item Het Bestuursorgaan beslist discretionair en zonder verdere motivatie dat een kandidaat niet wordt aanvaard als toegetreden lid.
    \item Het Bestuursorgaan kan met gewone meerderheid van de stemmen zonder verdere motivatie het lidmaatschap van een toegetreden lid schorsen.
        Deze schorsing geldt tot de eerstvolgende Algemene Vergadering waar beslist kan worden over het al dan niet uitsluiten van het toegetreden lid. Deze procedure verloopt conform artikel \ref{ArtikelQuoraAV}, \ref{ParUitsluitingLid}.
    \end{enumerate}
    
    \section{Einde van lidmaatschap}
    \begin{enumerate}
    \item Het lidmaatschap als werkend lid van de vereniging neemt een einde door de verstrijking van de termijn van lidmaatschap, vrijwillig ontslag, uitsluiting of overlijden. 
    \item Een werkend lid kan vrijwillig ontslag geven, door het ontslag toe te lichten aan het Bestuursorgaan of een formeel schrijven te richten aan de voorzitter.
    \item Een werkend lid wordt uitgesloten conform artikel \ref{ArtikelQuoraAV}, \ref{ParUitsluitingLid}.
    \item Een werkend lid dat na de officiële deadline van de KU Leuven geen student van de Faculteit Ingenieurswetenschappen is en geen deel uitmaakt van het Auditcomité, wordt van rechtswege uitgesloten als werkend lid, respectievelijk gewoon lid van de vereniging.
    \item Een werkend lid dat ontslag neemt of wordt uitgesloten, alsook de erfgenamen van een overleden werkend lid, kunnen geen aanspraak maken op het bezit van de vereniging en kunnen betaalde bedragen niet terugvorderen.
    \end{enumerate}


% \part{Praesidiumverkiezingen}\label{TitelVerkiezingen}

%     \section{Verkiezingen}\label{ArtikelVerkiezingen}
%     \begin{enumerate}
%     \item Het Bestuursorgaan is verplicht elk jaar tussen februari en juni verkiezingen te organiseren, de ``praesidiumverkiezingen''\ genaamd. Op deze verkiezingen zijn alle studenten aan de Faculteit Ingenieurswetenschappen aan de KU Leuven stemgerechtigd.
%     \item De precieze modaliteiten van de verkiezingen moeten worden vastgelegd in een specifiek verkiezingsreglement. Het Bestuursorgaan keurt het verkiezingsreglement goed. Het Bestuursorgaan moet het verkiezingsreglement openbaar maken op de website van de vereniging. Een aankondiging van deze publicatie dient te verschijnen in een papieren publicatie die gratis ter beschikking van de studenten aan de Faculteit Ingenieurswetenschappen gesteld wordt. Vanaf de bekendmaking van VKK (verkiezingscontrolecommissie) t.e.m. de dag van de jaarlijkse verkiezing mag het reglement geen wijzigingen ondergaan.
%     \item De organisatie van de verkiezingen verloopt eveneens met het oog op de verkiezingen van de vertegenwoordigers van de studenten bij de KU Leuven. Daarbij respecteert het Bestuursorgaan het Participatiedecreet en het Participatiereglement van Studentenraad KU Leuven en LOKO.
%     \item Deelnemen aan de verkiezingen dient in ploegverband te gebeuren. Een deelnemende ploeg moet mimimum vijf leden tellen. De leden van de verkiesbare ploeg moeten studenten aan de Faculteit Ingenieurswetenschappen zijn. Ten laatste twee weken voor de verkiezingsdag moet het Bestuursorgaan op de website van de vereniging de lijst met deelnemende ploegen publiceren, met inbegrip van de namen van de leden van alle ploegen.
%     \item Het verkiezingsreglement dient in elk geval te specificeren op welke manier bepaald wordt welke ploeg de verkiezingen heeft gewonnen. Er kan slechts één ploeg de verkiezingen winnen.
%     \end{enumerate}
    
%     \section{Bekrachtiging van de uitslag van de verkiezingen}
%     \begin{enumerate}
%     \item Het Bestuursorgaan ziet er op toe dat de bepalingen van artikel \ref{ArtikelVerkiezingen} van deze statuten en van het volledige verkiezingsreglement worden nageleefd. Vervolgens moet het Bestuursorgaan, na controle van het bepaalde in artikel \ref{ArtikelVerkiezingen} en in het verkiezingsreglement, beslissen of de verkiezingen conform de statuten en het verkiezingsreglement hebben plaatsgevonden. De bekrachtiging van de verkiezingsuitslag houdt in dat het Bestuursorgaan aanvaardt dat de leden van de winnende ploeg op de laatste Algemene Vergadering van het werkingsjaar toetreden als werkend lid.
%     \item Indien zij deze verplichting verzuimt, komt het de dagelijks bestuurders die het komende academiejaar niet langer student aan de Faculteit Ingenieurswetenschappen zullen zijn toe een aantal personen aan te wijzen die als werkend lid tot de vereniging zullen worden toegelaten, rekening houdend met het bepaalde in artikel \ref{ArtikelWerkendeLeden}. Deze dagelijks bestuurders kunnen discretionair en zonder motivatie de begin- en einddatum van het lidmaatschap van die personen bepalen er rekening mee houdende dat het lidmaatschap maximaal 1 jaar kan duren.  Deze personen staan in voor de continuïteit van de werking van VTK en zijn verplicht om nieuwe verkiezingen te organiseren.
%     \end{enumerate}


\part{De Algemene Vergadering}

    \section{De Algemene Vergadering}
    \begin{enumerate}
    \item De Algemene Vergadering bestaat uit alle werkende leden.
    \item Alle werkende leden hebben gelijk stemrecht. Elk werkend lid heeft één stem. De toegetreden leden maken geen deel uit van de Algemene Vergadering.
        \item De werkende leden zijn de praesidiumleden, twee afgevaardigden van elke erkende Werkgroep zoals gedefinieerd in het Huishoudelijk Reglement en maximaal vijf afgevaardigden van het Auditcomité.
    \end{enumerate}
    
    \section{Bevoegdheden van de Algemene Vergadering}
    \begin{enumerate}
    \item Conform de WVV-wet beschikt de Algemene Vergadering over de volgende bevoegdheden:
        \begin{itemize}
        \item de wijzigingen van de statuten;
        \item de benoeming en afzetting van bestuurders;
        \item de benoeming en afzetting van de commissaris en het bepalen van diens bezoldiging;
        \item de kwijting aan de bestuurders en de commissarissen;
        \item de goedkeuring van de begroting en van de jaarrekening;
        \item de ontbinding van de vereniging;
        \item de uitsluiting van een lid;
        \item de omzetting van de vzw in een ivzw, in een erkende coöperatieve vennootschap sociale onderneming of in een coöperatieve vennootschap erkend als sociale onderneming;
        \item een inbreng om niet van een algemeenheid te doen of te aanvaarden;
        \item de bepaling van de bezoldiging van de bestuurders in geval een bezoldiging wordt toegekend
        \item de bepaling van de bezoldiging van de commissaris
        \item het instellen van de verenigingsvordering tegen de bestuurders en de commissarissen
        \item alle andere gevallen waarin de wet of deze statuten dat vereisen.
        \end{itemize}
    \end{enumerate}
    
    \section{Bijeenkomst van de Algemene Vergadering}\label{ArtikelBijeenkomstenAV}
    \begin{enumerate}
    \item De jaarlijkse bijeenkomsten van de Gewone Algemene Vergadering zullen gehouden worden op de plaats, de dag en het uur vermeld in de uitnodiging. Deze Gewone Algemene Vergadering wordt gehouden binnen de vijf maanden na afsluiting van het boekjaar.
    \item Bijzondere bijeenkomsten in de Buitengewone Algemene Vergadering kunnen worden bijeengeroepen door de voorzitter van het Bestuursorgaan of op verzoek van ten minste 1/5 van de werkende leden. De voorzitter van het Bestuursorgaan roept de Algemene Vergadering bijeen binnen 21 dagen na het verzoek tot bijeenroeping en de Algemene Vergadering wordt uiterlijk 40 dagen na het verzoek gehouden.
    \item De vergaderingen worden bijeengeroepen door de voorzitter van het Bestuursorgaan. Aan de uitnodiging wordt een ontwerp van een agenda toegevoegd waarin elk punt is opgenomen dat door ten minste één bestuurder of één werkend lid wordt voorgesteld.
    \item De uitnodiging wordt ten minste 15 dagen voorafgaand aan de Algemene Vergadering verstuurd naar alle werkende leden.
    \end{enumerate}
    
    \section{Bepalingen bij de Algemene Vergaderingen}
    \begin{enumerate}
    \item De Algemene Vergadering kan geldig beraadslagen, ongeacht het aantal aanwezigen of vertegenwoordigde leden, tenzij de wet of de statuten het dwingend anders bepalen. Beslissingen worden genomen bij gewone meerderheid van de stemmen van de aanwezige of vertegenwoordigde werkende leden, tenzij de WVV-wet of deze statuten uitdrukkelijk anders vereisen. Stemming gebeurt bij handopsteking of, indien ten minste 1/3 van de werkende leden die aanwezig of vertegenwoordigd zijn erom vraagt, door geheime stemming.
    \item De Algemene Vergadering wordt voorgezeten door de voorzitter van het Bestuursorgaan. Indien deze niet aanwezig kan zijn stelt hij een vervanger aan. Deze vervanger wordt vermeld in de uitnodiging. In het geval er geen vervanger aangesteld is wordt de vergadering voorgezeten door de ondervoorzitter. Indien de ondervoorzitter eveneens afwezig is, wordt de Algemene Vergadering voorgezeten door de oudst aanwezige dagelijks bestuurder. Indien er geen dagelijks bestuurders aanwezig zijn, beslist de Algemene Vergadering onder leiding van de oudst aanwezige bestuurder of de bijeenkomst wordt uitgesteld. Indien er geen bestuurders aanwezig zijn, wordt de AV uitgesteld.
    \item Van elke vergadering worden notulen opgesteld, ondertekend door de voorzitter en de secretaris van het Bestuursorgaan, en bewaard in een notulenregister dat ter inzage zal zijn van de werkende leden, die hun inzagerecht zullen uitoefenen overeenkomstig de in de WVV-wet vastgelegde modaliteiten. Belanghebbende derden hebben inzagerecht op het zeteladres, na een vooraf gemaakte afspraak.
    \end{enumerate}

    
    \section{Quora van de Algemene Vergaderingen}\label{ArtikelQuoraAV}
    \begin{enumerate}
    \item De wijziging van de statuten vereist een beraadslaging in een Algemene Vergadering, die voldoet aan een quorum van 2/3 van de werkende leden door aanwezigheid of vertegenwoordiging. Ingeval op de eerste vergadering minder dan 2/3 van de leden aanwezig of vertegenwoordigd zijn, kan een tweede vergadering worden bijeengeroepen die geldig kan beraadslagen en besluiten alsook wijzigingen aannemen met de meerderheden hierna bepaald, ongeacht het aantal aanwezige of vertegenwoordigde werkende leden. De tweede vergadering mag niet binnen 15 dagen volgend op de eerste vergadering worden gehouden. De beslissing wordt geacht aanvaard te zijn indien ze wordt goedgekeurd door 2/3 van de stemmen van de aanwezige of vertegenwoordigde leden. Enkel wanneer de statutenwijziging betrekking heeft op het doel of het voorwerp waarvoor de vereniging is opgericht, vereist ze een meerderheid van 4/5 van de stemmen van de aanwezige of vertegenwoordigde werkende leden. \label{ParStatWijziging}
    \item De Algemene Vergadering kan besluiten tot de uitsluiting van een werkend of toegetreden lid met een meerderheid van 2/3 van de aanwezige of vertegenwoordigde werkende leden. Op de vergadering die over de uitsluiting beraadslaagt en beslist dient 2/3 van de werkende leden aanwezig of vertegenwoordigd te zijn. Indien de Algemene Vergadering beraadslaagt over het uitsluiten van een werkend of toegetreden lid, heeft dat lid in elk geval het recht om gehoord te worden. Indien het aanwezigheidsquorum niet gehaald wordt of de persoon die voorwerp is van de uitsluiting niet aanwezig kan zijn, kan een tweede AV samengeroepen worden zonder aanwezigheidsvereiste volgens dezelfde modaliteiten als \ref{ArtikelQuoraAV} \ref{ParStatWijziging} en waarbij het recht om gehoord te worden vervalt.
    \label{ParUitsluitingLid}
    \item De werkende leden kunnen zich op de vergadering laten vertegenwoordigen door één en slechts één ander werkend lid. Elk werkend lid kan drager zijn van slechts één volmacht. Personen die geen werkend lid zijn kunnen geen volmacht dragen.
    \end{enumerate}
    
    \section{Waarnemers}
    \begin{enumerate}
    \item Waarnemers kunnen de Algemene Vergadering bijwonen en mogen zich, mits toestemming van de voorzitter, tot de Algemene Vergadering richten.
    \item De Algemene Vergadering kan beslissen de gehele, of een deel van, de vergadering gesloten te verklaren.
    \item Indien de Algemene Vergadering gesloten wordt verklaard, zijn geen waarnemers meer toegelaten op de Algemene Vergadering, tenzij de Algemene Vergadering hier anders over beslist. Indien een deel van de Algemene Vergadering gesloten wordt verklaard, zijn geen waarnemers meer toegelaten op de betreffende delen van de Algemene Vergadering, tenzij de Algemene Vergadering hier anders over beslist.
    \end{enumerate}
    
    \section{Auditcomité}\label{ArtikelAudit}
    \begin{enumerate}
    \item De vereniging heeft een Auditcomité. Het Auditcomité brengt o.a. adviezen uit aan de Algemene Vergadering, voordat deze tot beslissing kan overgaan. Adviezen over onderstaande documenten zijn vereist voor de Algemene Vergadering tot een stemming kan overgaan:
        \begin{itemize}
        \item de begroting
        \item de jaarrekening
        \item kwijting
        \end{itemize}
    Indien een advies vereist is, maar dit niet werd uitgebracht, kan de Algemene Vergadering niet overgaan tot de bespreking en de stemming van het agendapunt tenzij ze met 2/3 meerderheid beslissen het agendapunt alsnog te bespreken en tot stemming over te gaan.
    \item De Algemene Vergadering kan te allen tijde bepalen om een advies aan het Auditcomité te vragen ter aanvulling van de punten van §1. Het Bestuursorgaan kan ook vrijwillig een advies vragen aan het Auditcomité en dit voorleggen aan de Algemene Vergadering.
    % \item De Algemene Vergadering kan, in het geval een advies vereist is, maar niet uitgebracht werd, niet overgaan tot de bespreking en de stemming van het agendapunt tenzij ze met 2/3 meerderheid beslissen om het agendapunt alsnog te bespreken en tot stemming over te gaan.
    \item Het Huishoudelijk Reglement neemt nadere bepalingen op omtrent het doel, de taken en de samenstelling van het Auditcomité; de termijnen en vormvereisten voor de adviezen en andere aspecten aangaande de werking van het Auditcomité.
    \end{enumerate}


\part{Het Bestuursorgaan}

    \section{Bestuurders}
    \begin{enumerate}
    \item De vereniging wordt bestuurd door een Bestuursorgaan, samengesteld uit ten minste vier personen. 
    \item \textcolor{green}{De vereniging wordt in alle rechtshandelingen en jegens derden geldig vertegenwoordigd door het gezamenlijk optreden van minstens twee bestuurders.}
    \item De bestuurders worden benoemd door een Algemene Vergadering, bij een gewone meerderheid van de stemmen van de aanwezige of vertegenwoordigde leden voor een termijn van één jaar of tot de eerstvolgende 15 juli. De bestuurders moeten op het moment van hun aanstelling werkend lid zijn van de vereniging.
    % \item Het Bestuursorgaan kiest onder de bestuurders een voorzitter (indachtig de voorwaarden in \ref{ParVoorwaardeVoorzitter} van dit artikel), ondervoorzitter, secretaris en minstens één penningmeester.
    % \item Het Bestuursorgaan zal bij het kiezen van een voorzitter rekening houden met volgende voorwaarde:
    %     \begin{itemize}
    %     \item De voorzitter moet in het jaar dat zijn mandaat van kracht is, zijn bachelordiploma aan de Faculteit Ingenieurswetenschappen behaald hebben en ingeschreven zijn voor een masteropleiding aan de Faculteit Ingenieurswetenschappen.
    %     \end{itemize}\label{ParVoorwaardeVoorzitter}
    \item Indien een persoon door de Algemene Vergadering tot bestuurder van de vereniging wordt benoemd, bepaalt de Algemene Vergadering de datum waarop dit mandaat ingaat. Deze datum is ten laatste de eerstvolgende 15 juli. Hun mandaat geldt tot de eerstvolgende 15 juli, te rekenen vanaf het moment dat hun mandaat ingaat. Het mandaat is onbeperkt hernieuwbaar.
    % \item De bestuurders voeren hun mandaat onbezoldigd uit.
    \item De benoeming van de leden van het Bestuursorgaan en hun ambtsbeëindiging worden openbaar gemaakt door neerlegging in het verenigingsdossier ter griffie van de ondernemingsrechtbank, en van een uittreksel daarvan bestemd om in de Bijlagen bij het Belgisch Staatsblad te worden bekendgemaakt. Uit de stukken moet in elk geval blijken of de personen die de vereniging vertegenwoordigen, de vereniging ieder afzonderlijk, gezamenlijk, dan wel als college verbinden alsook de omvang van hun bevoegdheden.
    \end{enumerate}
    
    \section{Einde van het bestuurderschap}
    \begin{enumerate}
    \item De bestuurders kunnen te allen tijde afgezet worden door de Algemene Vergadering die daarover beslist bij gewone meerderheid van de stemmen van de aanwezige of vertegenwoordigde werkende leden. 
    \item Bestuurders kunnen ook zelf ontslag nemen door een ontslagbrief, getekend door de bestuurder zelf, te richten aan het Bestuursorgaan of aan de voorzitter van het Bestuursorgaan. 
    \item Een bestuurder die door de omstandigheden geen werkend lid van de vereniging meer kan zijn, wordt van rechtswegen afgezet als bestuurder van de vereniging.
    \item Een bestuurder is verplicht om na zijn ontslag zijn opdracht verder te vervullen totdat redelijkerwijs in zijn vervanging kan worden voorzien voor een maximum van 1 maand. 
    \end{enumerate}

    \section{Praesidiumverkiezingen}
    \begin{enumerate}
        \item Het Bestuursorgaan is verplicht elk jaar verkiezingen te organiseren, de "praesidiumverkiezingen"\ genaamd. Op deze verkiezingen zijn alle studenten van de Faculteit Ingenieurswetenschappen aan de KU Leuven stemgerechtigd.
        \item De precieze modaliteiten van de praesidiumverkiezingen moeten worden vastgelegd in een specifiek verkiezingsreglement, dat door het Bestuursorgaan dient goedgekeurd te worden, en openbaar gemaakt moet worden op de website van de vereniging. Het reglement mag geen wijzigingen ondergaan vanaf de bekendmaking van VKK (verkiezingscontrolecommissie) t.e.m. de dag van de jaarlijkse verkiezingen. 
        \item De organisatie van de praesidiumverkiezingen verloopt eveneens met het oog op de verkiezingen van de vertegenwoordigers van de studenten bij de KU Leuven. Daarbij respecteerd het Bestuursorgaan het Participatiedecreet en het Participatiereglement van Studentenraad KU Leuven en LOKO.
        \item Het Bestuursorgaan moet beslissen of de verkiezingen conform de statuten en het verkiezingsreglement hebben plaatsgevonden. De bekrachtiging van de verkiezingsuitslag houdt in dat het Bestuursorgaan aanvaardt dat de leden van de winnende ploeg op de laatste Algemene Vergadering van het werkingsjaar toetreden als werkend lid. 
        \item Indien zij deze verplichting verzuimt, komt het de dagelijks bestuurders die het komend jaar niet meer studeren aan de Faculteit Ingenieurswetenschappen toe een aantal personen aan te wijzen als werkend lid van de vereniging. Deze personen staan in voor de continu\"iteit van de vereniging en zijn verplicht om nieuwe verkiezingen te organiseren. De dagelijks bestuurders bepalen discretionair en zonder motivatie de begin- en einddatum van het lidmaatschap van die personen, er rekening mee houdende dat het lidmaatschap maximaal \'e\'en jaar kan duren. 
    \end{enumerate}
    
    % \section{Bijeenkomst van het Bestuursorgaan}
    % \begin{enumerate}
    % \item Het Bestuursorgaan vergadert na oproeping door de voorzitter zo dikwijls als het belang van de vereniging het vereist. De vergadering wordt gehouden op de plaats die in het oproepingsbericht wordt vermeld.
    % \item De oproeping dient minstens 24 uur voor aanvang van de vergadering te gebeuren.
    % \item Het Bestuursorgaan wordt voorgezeten door de voorzitter. Indien deze niet aanwezig kan zijn stelt hij een vervanger aan. Deze vervanger wordt vermeld in de uitnodiging. In het geval er geen vervanger aangesteld is wordt het Bestuursorgaan voorgezeten door de ondervoorzitter. Indien de ondervoorzitter eveneens afwezig is, wordt het Bestuursorgaan voorgezeten door de oudst aanwezige dagelijks bestuurder. Indien deze allen afwezig zijn, beslist het Bestuursorgaan onder leiding van de oudst aanwezige bestuurder of de bijeenkomst wordt uitgesteld.
    % \item Het Bestuursorgaan kan slechts geldig vergaderen wanneer ten minste 1/4 van zijn leden aanwezig zijn. Het Bestuursorgaan oefent zijn bevoegdheden uit als college.
    % \item Van elke vergadering worden notulen opgesteld, ondertekend door de voorzitter en de secretaris van het Bestuursorgaan. De notulen kunnen tevens getekend worden door de bestuurders die het wensen. Ze worden bewaard in een notulenregister dat ter inzage zal zijn van de werkende leden, die hun inzagerecht zullen uitoefenen overeenkomstig de in artikel 3:103 van de WVV-wet vastgelegde modaliteiten.
    % \end{enumerate}
    
    % \section{Belangenvermenging}
    % \begin{enumerate}
    % \item Indien een bestuurder, rechtstreeks of onrechtstreeks, voor een beslissing die onder de bevoegdheid valt van het Bestuursorgaan een belang van vermogensrechtelijke of persoonlijke aard heeft dat strijdig is met het belang van de vereniging, moet hij dit meedelen aan de andere bestuurders voor het Bestuursorgaan een besluit neemt. De verklaring van deze bestuurder over de aard van het strijdig belang wordt opgenomen in de notulen van de vergadering.
    % \item Het Bestuursorgaan kan met gewone meerderheid van de stemmen beslissen of iemand al dan niet een belangenconflict van vermogensrechtelijke of persoonlijke aard heeft.
    % \item De bestuurder met het tegenstrijdig belang verwijdert zich uit de vergadering en onthoudt zich van de beraadslaging en de stemming over de aangelegenheid waarop het betrekking heeft.
    % \item Voormelde procedure is niet toepasselijk op gebruikelijke verrichtingen die plaatshebben onder de voorwaarden en tegen de zekerheden die op de markt gewoonlijk gelden voor soortgelijke verrichtingen.
    % \item Indien een bestuurder lid is van een (niet-)verkiesbare ploeg verwijdert hij zich uit de vergadering indien er een beraadslaging plaatsvindt die betrekking heeft op de controle op de verkiezingen of de bespreking van het verkiezingsresultaat.
    % \end{enumerate}
    
    % \section{Bevoegdheden van het Bestuursorgaan}
    % \begin{enumerate}
    % \item Het Bestuursorgaan is bevoegd om alle handelingen van intern bestuur te verrichten die nodig of dienstig zijn tot verwezenlijking van het doel van de vereniging, met uitzondering van die handelingen waarvoor volgens de wet of deze statuten de Algemene Vergadering exclusief bevoegd is.
    % \item Het Bestuursorgaan beslist over het Intern Reglement bij gewone meerderheid van de stemmen van de aanwezige of vertegenwoordigde werkende leden.
    % \item Onverminderd de verplichtingen die voortvloeien uit collegiaal bestuur, met name overleg en toezicht, kunnen de bestuurders de bestuurstaken onder elkaar verdelen. Zodanige verdeling van taken kan aan derden niet worden tegengeworpen, zelfs niet nadat ze zijn openbaar gemaakt. Niet-naleving ervan brengt wel de interne aansprakelijkheid van de betrokken bestuurder(s) in het gedrang.
    % \item Het Bestuursorgaan kan een deel van zijn bestuursbevoegdheid delegeren aan één of meerdere derden niet-bestuurders, zonder dat deze overdracht evenwel betrekking kan hebben op het algemeen beleid van de vereniging of de algemene bestuursbevoegdheid van het Bestuursorgaan. Hierin wordt een onderscheid gemaakt tussen vertegenwoordigers van werkgroepen en bijzondere mandaten. De bevoegdheid van vertegenwoordigers van werkgroepen wordt beschreven in het Intern Reglement, die van andere mandaten wordt precies afgebakend door het Bestuursorgaan, die ook de duur van het mandaat bepaalt. Het mandaat kan te allen tijde met onmiddellijke ingang door het Bestuursorgaan worden ingetrokken.
    % \end{enumerate}

    % \section{De plichten van het Bestuursorgaan}
    % \begin{enumerate}
    % \item De voorzitter van het Bestuursorgaan nodigt de leden van de Kringvergadering, bestaande uit alle verkozen praesidiumleden in een bepaald werkingsjaar zoals bepaald in Titel III, uit als waarnemer.
    % \item De voorzitter van het Bestuursorgaan stuurt de verslagen va het Bestuursorgaan door naar de leden van de Kringvergadering.
    % \item Het Bestuursorgaan bekrachtigd alle beslissingen van de Kringvergadering. Indien het Bestuursorgaan een beslissing van de ploegvergadering niet bekrachtigd, dient de voorzitter van het Bestuursorgaan alle leden van de Kringvergadering per mail op de hoogte te stellen van de genomen beslissing en deze beslissing te motiveren.
    % \end{enumerate}
    
    
    % \section{Vertegenwoordiging}\label{ArtikelVertegenwoordiging}
    % \begin{enumerate}
    % \item Het Bestuursorgaan vertegenwoordigt als college de vereniging in alle handelingen in en buiten rechte.
    % \item Onverminderd de algemene vertegenwoordigheidsbevoegdheid van het Bestuursorgaan als college, wordt de vereniging in en buiten rechte eveneens geldig vertegenwoordigd door ofwel:
    %     \begin{itemize}
    %     \item twee dagelijks bestuurders, waarvan één van beide voorzitter, ondervoorzitter, penningmeester of secretaris is; ofwel
    %     \item één dagelijks bestuurder en één bestuurder tot een bedrag van maximaal EUR 7500,-; ofwel
    %     \item één dagelijks bestuurder tot een bedrag van EUR 2500,- voor wat het dagelijks bestuur betreft dat binnen de bevoegdheid valt van deze persoon zoals beschreven in het Intern Reglement.
    %     \end{itemize}\label{ParLimietenVertegenwoordiging}
    % \item Het Bestuursorgaan kan, wanneer de noodzakelijkheid en het belang van de vereniging zulks vereisen, personen uitzonderlijk machtiging geven om ten aanzien van derden te handelen tot een door het Bestuursorgaan gespecificeerd bedrag.
    % \end{enumerate}
    
    \section{Dagelijks bestuur}
    \begin{enumerate}
    \item Het dagelijks bestuur van de vereniging alsook de vertegenwoordiging wat betreft dat bestuur kan door het Bestuursorgaan worden opgedragen aan één of meer van haar leden.
    \item De dagelijks bestuurder(s) is (zijn) wat betreft externe vertegenwoordigingsmacht bevoegd te handelen volgens de bepalingen in het Huishoudelijk Reglement
    \item Het dagelijks bestuur omvat zowel de handelingen en de beslissingen die niet verder reiken dan de behoeften van het dagelijks leven van de vereniging, als de handelingen en de beslissingen die, ofwel om reden van hun minder belang dat ze vertonen, ofwel omwille van hun spoedeisend karakter, de tussenkomst van het Bestuursorgaan niet rechtvaardigen.
    \item De benoeming van de personen belast met het dagelijks bestuur en hun ambtsbeëindiging wordt openbaar gemaakt door neerlegging in het verenigingsdossier ter griffie van de ondernemingsrechtbank, en van een uittreksel daarvan bestemd om in de Bijlagen bij het Belgisch Staatsblad te worden bekendgemaakt. Uit de stukken moet in ieder geval blijken of de personen die vereniging vertegenwoordigen inzake dagelijks bestuur, de vereniging ieder afzonderlijk, gezamenlijk, dan wel als college verbinden alsook de omvang van hun bevoegdheden.
    \item In uitzonderlijke gevallen, wanneer de dringende noodzakelijkheid en het belang van de vereniging zulks vereisen, kunnen de besluiten van het Bestuursorgaan worden genomen bij eenparig schriftelijk akkoord van de dagelijks bestuurders. Daartoe is vereist dat er vooraf een unaniem akkoord is onder de dagelijks bestuurders om tot schriftelijke besluitvorming over te gaan. Schriftelijke besluitvorming veronderstelt in elk geval dat er een beraadslaging plaatsvond per e-mail, video- of telecommunicatie.
    \end{enumerate}
    
    \section{Aansprakelijkheid}
    \begin{enumerate}
    \item De bestuurders en dagelijks bestuurders zijn niet persoonlijk verbonden door de verbintenissen van de vereniging. Tegenover de vereniging en tegenover derden is hun verantwoordelijkheid beperkt tot de vervullingen van de hun gegeven opdracht. Overeenkomstig het gemeen recht, het bepaalde in de wet en deze statuten zijn ze aansprakelijk voor de tekortkomingen in hun bestuur.
    \end{enumerate}


\part{Diverse bepalingen}

    % \section{Volmachten}
    % \begin{enumerate}
    % \item De penningmeester(s) krijgen een volmacht op de rekeningen van de vereniging. De penningmeester(s) hebben een volmacht om, conform het Intern Reglement, volmachten voor specifieke rekeningen van de vereniging uit te delen aan andere bestuurders en de werkgroepen.
    % \item De voorzitter van de vereniging krijgt een volmacht op de rekening(en) van de vereniging. 
    % \end{enumerate}
    
    \section{Boekhouding}
    \begin{enumerate}
    \item Het boek- en werkjaar begint op 1 augustus en eindigt op 31 juli van het jaar erna. In uitzonderlijke gevallen kan gebruik gemaakt worden van een verlengd boekjaar. De toestemming hiervoor wordt gegeven door de Algemene Vergadering.
    \item De boekhouding wordt gevoerd conform de WVV-wet.
    \item De jaarrekening wordt neergelegd in het dossier gehouden ter griffie van de ondernemingsrechtbank overeenkomstig het bepaalde in de WVV-wet. Voor zover van toepassing wordt de jaarrekening neergelegd bij de Nationale Bank overeenkomstig het bepaalde in de WVV-wet.
    \item Het Bestuursorgaan legt het jaarresultaat van het voorafgaande boekjaar voor aan de Algemene Vergadering. Bij goedkeuring van deze documenten kunnen de bestuurders gekwijt worden.
    \item Het Bestuursorgaan legt een voorstel van de begroting van het boekjaar voor aan de Algemene Vergadering. 
    \end{enumerate}
    
    \section{De ontbinding}
    \label{ArtikelOntbinding}
    \begin{enumerate}
    \item De Algemene Vergadering zal worden samengeroepen ter bespreking van voorstellen m.b.t. de ontbinding voorgelegd door het Bestuursorgaan of door minimum 1/5 van alle werkende leden. De samenroeping en agendering vinden plaats overeenkomstig het bepaalde in artikel \ref{ArtikelBijeenkomstenAV} van deze statuten.
    \item De beraadslaging en beslissing over de ontbinding respecteert het quorum en de meerderheid in artikel \ref{ArtikelQuoraAV} van deze statuten. Dezelfde regels gelden (inzake aanwezigheid/vertegenwoordiging en stemming) als voor de wijziging van het doel of het voorwerp. Vanaf de beslissing tot ontbinding vermeldt de vereniging steeds dat zij "vzw in vereffening"\ is overeenkomstig artikel 2:115 van de WVV-wet.
    \item Ingeval het voorstel tot ontbinding wordt goedgekeurd, benoemt de Algemene Vergadering één of meerdere vereffenaar(s) waarvan zij de opdracht zullen omschrijven.
    \item Ingeval van ontbinding en vereffening, beslist de Algemene Vergadering over de bestemming van het vermogen van de vereniging dat moet worden toegekend aan een andere vereniging zonder winstoogmerk met een gelijkaardig of verwant doel.
    \item Alle beslissingen betreffende de ontbinding, de vereffeningsvoorwaarden, de benoeming en de ambtsbeëindiging van de vereffenaars, de afsluiting van de vereffening en de bestemming van het actief worden neergelegd ter griffie van de ondernemingsrechtbank en bekendgemaakt in de Bijlage bij het Belgisch Staatsblad overeenkomstig het bepaalde in de WVV-wet.
    \end{enumerate}
    
    \section{Aanvullende bepalingen}
    \begin{enumerate}
    \item Voor alles wat door deze statuten niet uitdrukkelijk is geregeld, is de WVV-wet van toepassing. De Algemene Vergadering moet een Huishoudelijk Reglement ter beschikking stellen ter toelichting en uitbreiding van deze statuten. Het Bestuursorgaan moet een Intern Reglement ter beschikking stellen ter toelichting en uitbreiding van deze statuten.
    \end{enumerate}
    
    \section{Historiek der wijzigingen}
    \begin{enumerate}
    \item Deze geco\"ordineerde statuten werden aangenomen op de AV van 't ElixIr vzw op 3 november 2002.
    \item Deze geco\"ordineerde statuten werden gewijzigd op de AV van 't ElixIr vzw op 1 juli 2005.
    \item Deze geco\"ordineerde statuten werden gewijzigd op de AV van 't ElixIr vzw op 12 november 2006.
    \item Deze geco\"ordineerde statuten werden gewijzigd op de AV van 't ElixIr vzw op 2 september 2007.
    \item Deze geco\"ordineerde statuten werden gewijzigd op de AV van 't ElixIr vzw op 26 augustus 2008.
    \item Deze geco\"ordineerde statuten werden gewijzigd op de AV van VTK Ondersteuning vzw op 16 november 2008.
    \item Deze geco\"ordineerde statuten werden gewijzigd op de AV van VTK Ondersteuning vzw op 28 november 2010.
    \item Deze geco\"ordineerde statuten werden gewijzigd op de AV van VTK Ondersteuning vzw op 22 oktober 2012.
    \item Deze geco\"ordineerde statuten werden gewijzigd op de AV van VTK Ondersteuning vzw op 8 november 2015.
    \item Deze geco\"ordineerde statuten werden gewijzigd op de AV van VTK Ondersteuning vzw op 1 mei 2016.
    \item Deze geco\"ordineerde statuten werden gewijzigd op de AV van VTK Ondersteuning vzw op 19 mei 2021.
    \item Deze geco\"ordineerde statuten werden gewijzigd op de AV van VTK Ondersteuning vzw op 1 juli 2021.
    \item Deze geco\"ordineerde statuten werden gewijzigd op de AV van VTK Ondersteuning vzw op 14 juli 2022. 
    \item Deze geco\"ordineerde statuten werden gewijzigd op de AV van Vlaamse Technische Kring vzw op 3 juli 2023.
    \item Deze geco\"ordineerde statuten werden gewijzigd op de AV van Vlaamse Technische Kring vzw op 1 juli 2024.
    \item Deze geco\"ordineerde statuten werden gewijzigd op de AV van Vlaamse Technische Kring vzw op 6 juli 2026.
    \end{enumerate}

\end{document}